%Respectau títol del capítol.
\chapter{Acrònims}
\label{ch:acronims}
%
% Per utilitzar els acrònims es recomana fer un poc
% de recerca bibliogràfica per entendre com
% funcionen. Concretament podeu llegir el manual
% que teniu dins el vostre sistema.
% La comanda `texdoc acronym` hauria de mostrar-lo.
%
\begin{acronym}

    \acro{ABI}[ABI]{Application Binary Interface}
    \acro{ADR}[ADR]{Architecture Decision Record}
    \acro{AMM}[AMM]{Automated Market Maker}
    \acro{API}[API]{Application Programming Interface}
    \acro{ARC}[ARC]{Algorand Request for Comments}
    \acro{AVM}[AVM]{Algorand Virtual Machine}
    \acro{CI}[CI]{Continuous Integration}
    \acro{CSV}[CSV]{Comma-Separated Values}
    \acro{DoD}[DoD]{Definition of Done}
    \acro{DoR}[DoR]{Definition of Ready}
    \acro{EVM}[EVM]{Ethereum Virtual Machine}
    \acro{IdP}[IdP]{Identity Provider}
    \acro{JSON}[JSON]{JavaScript Object Notation}
    \acro{KMD}[KMD]{Key Management Daemon}
    \acro{MVP}[MVP]{Minimum Viable Product}
    \acro{PR}[PR]{Pull Request}
    \acro{RFC}[RFC]{Request for Comments}
    \acro{RPC}[RPC]{Remote Procedure Call}
    \acro{SHA}[SHA]{Secure Hash Algorithm}
    \acro{TEAL}[TEAL]{Transaction Execution Approval Language}
    \acro{TFG}[TFG]{Treball de Fi de Grau}
    \acro{UIB}[UIB]{Universitat de les Illes Balears}
    \acro{US}[US]{User Story}
    \acro{ZK}[ZK]{Zero-Knowledge}
    \acro{ZK-SNARK}[ZK-SNARK]{Zero-Knowledge Succinct Non-Interactive Argument of Knowledge}

\end{acronym}
